\chapter*{Conclusion générale et perspectives}
\addcontentsline{toc}{chapter}{Conclusion générale et perspectives}
\markboth{Conclusion générale et perspectives}{Conclusion générale et perspectives}

\section*{Réponse à la problématique d'ingénierie}

Le travail mené tout au long de ce projet de fin d'études a permis d'apporter une réponse concrète, opérationnelle et scientifiquement rigoureuse à la problématique de la digitalisation sécurisée et économique de la gestion scolaire en Afrique de l'Ouest.

Face à la prolifération de logiciels monopostes vulnérables et au coût prohibitif des solutions propriétaires occidentales, la plateforme \textbf{École~221} démontre la faisabilité d'un modèle SaaS multi-tenant mutualisé capable de diviser par dix le coût d'accès au numérique pour les établissements scolaires sénégalais. Grâce au découpage hexagonal et à l'intercepteur de sécurité cryptographique (\texttt{SEC-003}), l'étanchéité absolue des données entre écoles est garantie au niveau applicatif et persistée en base sans aucune dépendance envers un paramètre manipulable côté client.

\section*{Bilan des objectifs et validation par les preuves}

Tous les objectifs fonctionnels et de sécurité initialement assignés au prototype ont été atteints et formellement validés par automate :

\begin{itemize}
  \item \textbf{16 modules fonctionnels} ont été réalisés et qualifiés de bout en bout ;
  \item \textbf{142 assertions de test d'intégration} exécutées avec un taux de conformité de \textbf{100\% (PASS)} consigné dans le journal d'exécution (\nolinkurl{logs/qa_execution.log}) ;
  \item \textbf{Étanchéité multi-tenant prouvée} par les tests de sécurité de l'Agent~07 interdisant tout contournement d'accès ;
  \item \textbf{Résilience hors-ligne opérationnelle} pour les agents vigiles aux portes du campus, garantissant le contrôle d'accès même en cas de coupure du réseau Internet ;
  \item \textbf{Chiffrage économique viable} démontrant un seuil de rentabilité (\emph{point mort}) atteint dès 2 écoles partenaires pour un coût cloud moyen de seulement 29,38~USD par mois.
\end{itemize}

\section*{Limites constatées du travail réalisé}

Dans un esprit de parfaite honnêteté intellectuelle et académique, deux limites techniques majeures du prototype actuel doivent être explicitées :

\begin{enumerate}
  \item \textbf{L'absence d'intégration directe d'une passerelle de paiement Mobile Money en production} : si l'architecture prévoit les interfaces d'encaissement et simule les notifications de paiement, le branchement effectif avec les APIs d'opérateurs locaux (Wave, Orange Money) nécessite l'obtention d'un agrément d'opérateur marchand qui dépasse le cadre académique de ce mémoire ;
  \item \textbf{L'absence de qualification en charge extrême} : les tests de performance ont validé les temps de réponse nominaux sous une charge de 500 requêtes concurrentes, mais le comportement de la base de données au-delà de 50~000 élèves simultanés reste à qualifier sur un environnement de banc d'essai dédié.
\end{enumerate}

\section*{Perspectives d'évolution et recherche future}

Les perspectives ouvertes par ce projet d'ingénierie sont nombreuses et structurantes :

\begin{itemize}
  \item \textbf{Déploiement d'une application mobile native pour les parents} : concevoir avec React Native une déclinaison mobile optimisée pour les notifications push en temps réel des absences et notes ;
  \item \textbf{Moteur d'alerte précoce contre le décrochage scolaire} : intégrer un module d'apprentissage automatique supervisé capable d'alerter les conseillers pédagogiques dès qu'un élève présente une chute anormale de moyenne ou un taux d'absentéisme croissant ;
  \item \textbf{Généralisation du framework agentique} : capitaliser sur l'expérience acquise avec le pipeline multi-agents pour standardiser le développement de futures applications publiques en Afrique de l'Ouest.
\end{itemize}
